% Section 7 — Out-of-Sample: The Pre-Registered Tourism Test
% Numbers sourced from: paper_rewrite_plan.md §6 (§6 subplan, Tourism bullet), §19
% (corrections table, trend-branch row), §20 (Tourism preset validation row);
% response_AC.md Point 1; response_ovPT.md "Out of sample".

\section{Out-of-Sample: The Pre-Registered Tourism Test}\label{sec:tourism}

The findings of the preceding sections condense to a single compact configuration,
fixed in writing before the test began: ten stacks, a 128-wide trunk, a frozen
complete-rank forecast basis ($\df = \Ho$) of any family, a truncated backcast basis
($\db < \L$), a rectified output head, and a paired trend branch. No element of this
configuration was adjusted on the test data. The test bed is the Tourism competition
set: 1{,}311 series across the Yearly, Quarterly, and Monthly horizons, evaluated over
240 runs and three horizons. Because the configuration and the acceptance criteria
predate the runs, the test distinguishes what the preceding measurements actually
predict from what they merely describe.

Four predictions were registered alongside the configuration.
Table~\ref{tab:tourism_predictions} lists each with its registered value and its
outcome. Two of four registered predictions failed and are reported as failures.

\begin{table}[t]
  \centering
  \caption{Registered predictions for the pre-registered Tourism out-of-sample test
    and their outcomes. The registered acceptance threshold for P1 was 0.5 \smape{}
    against the better of the two reference baselines, N-BEATS-G and N-BEATS-I+G
    (Oreshkin et al., 2020). Two of the four predictions failed; both failures are
    reported as failures rather than rescoped.}
  \label{tab:tourism_predictions}
  \begin{tabular}{@{}p{0.26\linewidth}p{0.34\linewidth}p{0.34\linewidth}@{}}
    \toprule
    Prediction & Registered value & Outcome \\
    \midrule
    P1: the compact preset matches the better baseline
      & Within 0.5 \smape{} of the better baseline on all three horizons, at
        14--15$\times$ fewer parameters
      & \textbf{Confirmed} on all three horizons; zero collapses in 240 runs \\
    \addlinespace
    P2: the narrow trunk holds
      & $\U = 128$ against $\U = 512$
      & \textbf{Confirmed} on all three horizons; significant on Quarterly only
        (uncorrected $p = 0.031$) \\
    \addlinespace
    P3: the paired trend branch is required
      & Branch always on
      & \textbf{Failed}; demoted to conditional. The branch costs 0.344 \smape{} on
        Quarterly ($p = 0.021$) and 0.609 on Monthly ($p = 0.0073$); what it buys is
        tail protection at the shortest horizon, where at $\Ho = 4$ it cuts the
        per-seed standard deviation from 8.12 to 0.46 \\
    \addlinespace
    P4: parity between the in-block basis design and the sub-band design
      & Budget-matched parity
      & \textbf{Failed}; the in-block basis wins outright
        (Sec.~\ref{sec:scoped}) \\
    \bottomrule
  \end{tabular}
\end{table}

\paragraph{Interpretation.}
The trunk result of Sec.~\ref{sec:trunk} holds out of sample. A measurement taken on
M4 Yearly and Monthly --- trunk width inert far below the published defaults ---
predicted the behavior of a dataset family it had never seen, with no retuning: the
128-wide trunk matches or beats the 512-wide trunk on every horizon, and the one
significant cell favors the narrow trunk. The trend branch, by contrast, is
conditional rather than universal. On the two longer horizons it is a significant
cost, 0.3--0.6 \smape{}; at the shortest horizon it is the difference between a stable
configuration and an unstable one. The carried-forward rule is therefore keyed to the
conditions that predict the failure --- low effective rank or a dominant window mean
--- not applied unconditionally. The zero-collapse count is likewise read at its
registered scope: the rectified head removed the sign-failure mode of
Sec.~\ref{sec:trainability} on this non-negative target domain, which is the mode it
was designed to remove. All generality evidence in this section is scoped to
forecasting.
